\scanpage{199}
\mspage{8.}
On explicite facilmt la composition des hom des données de ramification associées à $Z,G$, $Z',G'$, $Z'',G''$ \ldots.. Ceci dit :

Soient des $S'_i$ couvrant $S$ ; sur chaque $S_i$ supposons donné un $(Z_i,G_i)$ schéma à groupe d'opérateurs $G_i$ admissible ; pour tt couple $(i,j)$, on suppose donné \deleted{\unsure{un}} sur $S_i\times_S S_j$ une ``\uline{subordination}'' de $(Z_{ij},G_i)$ à $(Z_{ji},G_j)$ $(Z_{ij}=Z_i\times_{S_i}(S_i\times_S S_j))$, \deleted{de telle façon qu'on ait} sur les $S_{ijk}$ des isomorphismes de transitivité entre les subordinations. On suppose enfin une condition de compatibilité sur les $S_{ijkl}$. Alors on peut recoller \deleted{les} $(\mathscr R_i,r_i)$ définis par les $(Z_i,G_i)$ \ldots..

Donnons un cas particulier moins enclin à donner le mal de mer :

Soit une catégorie fibrée \inserted{(de ``modèles'')} $C$ sur $Et(S)$ (pas néc. un champ) et un foncteur cartésien $s:C\unsure{\longrightarrow|-}\mathrm{Fl}_S$. On suppose $C$ fibrée en groupoïdes, que deux objets d'un $C_{S'}$ sont localmt isomorphes, que $C_S$ est localement non vide -- d'où finalement, en prenant son champ associé, une \uline{gerbe}. On suppose que le lien de cette gerbe est défini \deleted{\unsure{par un objet}} loc. constant tordu, i.e. pour tt $x\in C_{S'}$, le faisceau associé
% archive sheet number: 198
